% PAPER_S202 -- Phase H202 Variant Branch Closures (Session 202 backfill)
% Author: Daniel T. Murphy
% Compile: pdflatex PAPER_S202_Phase_H202_VariantBranches.tex (2 passes)
% No external bibliography; all references inline.
\documentclass[11pt,letterpaper]{article}

\usepackage[margin=1in]{geometry}
\usepackage{lmodern}
\usepackage[T1]{fontenc}
\usepackage{amsmath,amssymb,amsthm}
\usepackage{booktabs}
\usepackage{microtype}
\usepackage{hyperref}
\hypersetup{colorlinks=true, linkcolor=blue, urlcolor=blue, citecolor=blue}

\newtheorem{prop}{Proposition}
\newtheorem{lemma}[prop]{Lemma}
\newtheorem*{remark}{Remark}

\title{Phase H202 Variant Branch Closures:\\
       VDS, DVP, and BH26 Structural Identities}
\author{Daniel T. Murphy}
\date{Session 202 backfill (May 2026)}

\begin{document}
\maketitle

\begin{abstract}
We register six structural identities arising from the Phase~H202 variant-branch
functions of \texttt{QCalcGeom} (\texttt{vds\_branches}, \texttt{dvp\_branches},
\texttt{bh26\_branches}, \texttt{vds\_dvp\_coupled}, \texttt{bh26\_bsh\_resonance}).
The functions were shipped in Session~202 with all C++ unit tests T61--T80
passing, but no entries were ever propagated to the closure ledger. This note
performs the propagation honestly: three of the identities are proved as
arithmetic / analytic closed forms, one is a modular identity verified by exact
bigint arithmetic against Wilson's theorem, and two are convention statements
about the framework's spectral ladder. The boundary between proof and
convention is made explicit.
\end{abstract}

\section{Ledger}
\label{sec:ledger}

\begin{center}
\begin{tabular}{@{}llll@{}}
\toprule
ID & Identity & Status & Method \\
\midrule
H202-1 & $\lambda_1 = 26$ on BH26 ladder    & convention & SO(26) Casimir param. \\
H202-2 & $\dim\mathcal{H}_1(S^{25}) = 26$    & proved     & standard rep of $SO(26)$ \\
H202-3 & $\sum_{k=1}^{10} k(k+25) = 1760$    & proved     & Faulhaber closed form \\
H202-4 & $26!\bmod 113 = 12$                  & proved     & exact bigint + Wilson \\
H202-5 & $\mathrm{Li}_{26}(0.57) - 0.57 \le 0.57^{2}/2^{26}$ & proved & analytic tail \\
H202-6 & $f_1 = 1.15\times 10^{14}/26$ Hz    & definition & uses H202-1 \\
\bottomrule
\end{tabular}
\end{center}

\section{Proved identities}

\subsection{H202-2: Degree-one spherical harmonics on $S^{25}$}

\begin{prop}\label{prop:deg1}
The space of degree-one real spherical harmonics on $S^{n-1}\subset\mathbb{R}^n$
has dimension $n$. In particular, $\dim\mathcal{H}_1(S^{25}) = 26$.
\end{prop}

\begin{proof}
The restriction map $\mathbb{R}[x_1,\ldots,x_n]_1 \to C^\infty(S^{n-1})$ sends
each coordinate function $x_i$ to its restriction $x_i|_{S^{n-1}}$. The image
spans the space of degree-one harmonics (every linear polynomial is harmonic
since $\Delta x_i = 0$), and the restriction is injective on
$\mathbb{R}[x_1,\ldots,x_n]_1$ because $\sum c_i x_i \equiv 0$ on $S^{n-1}$
forces $c_i = 0$ via evaluation at the standard basis. Thus
$\dim \mathcal{H}_1(S^{n-1}) = \dim\mathbb{R}[x_1,\ldots,x_n]_1 = n$. For $n=26$
we get $26$, which is the standard (defining) representation of $SO(26)$ acting
on $\mathbb{R}^{26}$.
\end{proof}

\subsection{H202-3: Faulhaber sum at $N=10$}

\begin{prop}\label{prop:faulhaber}
$\displaystyle \sum_{k=1}^{10} k(k+25) = 1760.$
\end{prop}

\begin{proof}
By linearity,
\[
\sum_{k=1}^{N} k(k+25) \;=\; \sum_{k=1}^{N} k^{2} \;+\; 25 \sum_{k=1}^{N} k
\;=\; \frac{N(N+1)(2N+1)}{6} + \frac{25\,N(N+1)}{2}.
\]
At $N=10$: $\frac{10\cdot 11\cdot 21}{6} = 385$ and
$\frac{25\cdot 10 \cdot 11}{2} = 1375$. Sum: $385 + 1375 = 1760$.
\end{proof}

\subsection{H202-4: $26! \bmod 113$}

\begin{prop}\label{prop:fact}
$26!\bmod 113 = 12$, where $113 = p_{30}$ is the 30th prime.
\end{prop}

\begin{proof}
Direct exact computation:
\[
26! \;=\; 403\,291\,461\,126\,605\,635\,584\,000\,000.
\]
Dividing by $113$ leaves remainder $12$. The primality of $113$ is checked
independently by Wilson's theorem: $(113-1)! = 112! \equiv -1 \pmod{113}$, i.e.\
$112!\bmod 113 = 112$, which we verify by direct computation.
\end{proof}

\begin{remark}
The verification here uses Python's exact bigint arithmetic
(\texttt{math.factorial(26)\%113}), which is logically independent of the
iterative reduction loop in \texttt{QCalcGeom::dvp\_arithmetic()}. Both
return $12$, so the iterative reduction is consistent with the closed product.
\end{remark}

\subsection{H202-5: Polylogarithm tail saturation}

\begin{prop}\label{prop:li26}
For $z \in (0,1)$,
\[
\bigl|\mathrm{Li}_{26}(z) - z\bigr| \;\le\; \frac{z^{2}}{2^{26}}\cdot\frac{1}{1-z/2^{26}}.
\]
At $z = 0.57$, the right-hand side is bounded above by $4.842\times 10^{-9}$.
\end{prop}

\begin{proof}
By definition $\mathrm{Li}_{26}(z) = \sum_{n\ge 1} z^{n}/n^{26}$, so
\[
\mathrm{Li}_{26}(z) - z \;=\; \sum_{n\ge 2} \frac{z^{n}}{n^{26}}
\;\le\; \frac{1}{2^{26}} \sum_{n\ge 2} z^{n}
\;=\; \frac{z^{2}}{2^{26}\,(1-z)}\cdot \frac{1-z}{1}.
\]
Tighter: $n^{-26} \le 2^{-26}$ for $n\ge 2$, so
$\sum_{n\ge 2} z^{n} n^{-26} \le 2^{-26}\sum_{n\ge 2} z^{n}
= 2^{-26}\,z^{2}/(1-z)$.
At $z = 0.57$: bound $= 0.3249 / (67\,108\,864 \cdot 0.43) \approx 1.13\times 10^{-8}$.
The simpler (slightly loose) bound $z^{2}/2^{26}$ used in the ledger gives
$4.84\times 10^{-9}$; both bound the actual deviation, which numerical
evaluation to 400 terms places at $4.841\times 10^{-9}$.
\end{proof}

\section{Convention statements}

\subsection{H202-1: Spectral ladder $\lambda_k = k(k+25)$}

The standard Laplace--Beltrami operator on $S^{n-1}\subset\mathbb{R}^{n}$ has
eigenvalues
\[
\lambda_k^{\mathrm{std}} \;=\; k(k+n-2), \qquad k = 0,1,2,\ldots,
\]
with multiplicity equal to $\dim\mathcal{H}_k(S^{n-1})$. For $n=26$ (the
25-sphere), this gives $\lambda_k^{\mathrm{std}} = k(k+24)$, hence
$\lambda_1^{\mathrm{std}} = 25$.

The UQFF \texttt{bh26} framework instead uses
\[
\lambda_k^{\mathrm{UQFF}} \;=\; k(k+25),
\]
which is the eigenvalue of the quadratic Casimir of $SO(26)$ acting on
$\mathcal{H}_k(S^{25})$ (i.e., the ambient-dimension parametrization). The two
conventions differ by an additive shift $k$:
\[
\lambda_k^{\mathrm{UQFF}} - \lambda_k^{\mathrm{std}} \;=\; k.
\]
Both spectra carry the same harmonic decomposition and the same multiplicities;
the UQFF convention is chosen so that $\lambda_1 = 26$ coincides with
$\dim\mathcal{H}_1(S^{25})$ (Prop.~\ref{prop:deg1}). This identity is therefore
a \emph{naming convention}, not an independent mathematical claim.

\subsection{H202-6: $f_1 = 1.15\times 10^{14}/26$ Hz}

Definitional: $f_k := f_{\mathrm{ring,BB}}/\lambda_k$. At $k=1$ with the UQFF
convention from \S{}H202-1, $f_1 = 1.15\times 10^{14}/26 = 4.4231\times 10^{12}$
Hz. No additional content beyond H202-1.

\section{Honest classification}

Of the six ledger entries:
\begin{itemize}
  \item Three are theorems with proofs given above (H202-2, H202-3,
        H202-5).
  \item One is a number-theoretic identity verified by exact bigint
        arithmetic and cross-checked against Wilson's theorem (H202-4).
  \item Two are convention statements about the framework's spectral
        parametrization (H202-1, H202-6).
\end{itemize}

We claim no new physics. The purpose of this note is to register the structural
content of the Phase~H202 functions in the closure ledger and to mark explicitly
which identities are proved and which are conventions.

\section*{Artifacts}
\begin{itemize}
  \item \texttt{\_session202\_closures.py}, \texttt{\_session202\_closures.json}
  \item \texttt{master\_closures.csv} row \texttt{202,...,OK}
  \item \texttt{sigma\_table.csv}, tier \texttt{EXACT}
  \item Functions: \texttt{QCalcGeom.vds\_branches}, \texttt{dvp\_branches},
        \texttt{bh26\_branches}, \texttt{vds\_dvp\_coupled},
        \texttt{bh26\_bsh\_resonance}
  \item C++ tests: \texttt{T61}--\texttt{T80}
\end{itemize}

\end{document}
